\section{Incidencia hexada–octada y funcional orientado}
\subsection{Incidencia dodecafásica y funcional de retorno}

La contribución de Barbero--Immirzi utiliza el transporte angular y la
incidencia hexada--octada. Sean $\mathsf H\subset\mathsf O$ una hexada
y una octada con el par marcado de la hexada conservado. Los conjuntos de
banderas pertinentes son
\begin{equation}
 \mathcal F_6=\mathsf H\times\binom{\mathsf H}{2}
 \lhook\joinrel\longrightarrow
 \mathcal F_8=\mathsf O\times\binom{\mathsf H}{2},
 \qquad |\mathcal F_6|=90,\quad |\mathcal F_8|=120.
 \label{iii:eq:barbero-flags}
\end{equation}
Los cardinales son $6\binom62$ y $8\binom62$; la misma incidencia
conserva el defecto normalizado
$(90-120)/(24\binom62)=-1/12$. Esta cuenta identifica los dos sectores
empleados por el retorno, mientras el calendario determina sus
multiplicidades relativas.

Para $m\in\{90,120\}$, el retorno trítico sobre un canal contractivo es
\begin{equation}
 S_m(q)=\sum_{j\geq0}q^{m+3mj}
       =\frac{q^m}{1-q^{3m}},\qquad 0<q<1.
 \label{iii:eq:barbero-return-series}
\end{equation}
Cada término corresponde a la altura inicial $m$ y a $j$ retornos de
longitud $3m$. La serie geométrica conserva, por tanto, tanto la altura
inicial como el período del retorno. El calendario dodecafásico
$\mathbb Z/12\mathbb Z$ aporta doce sectores y una única incidencia
orientada de frontera, denotada $\partial_{\rm ret}$. En el módulo
libre de sucesos, su cadena fundamental es
\begin{equation}
 c_{12}^{+}
 =\sum_{j\in\mathbb Z/12\mathbb Z}[j,\mathcal F_6]
  -[\partial_{\rm ret},\mathcal F_8].
 \label{iii:eq:barbero-cycle}
\end{equation}
El lector lineal $\mathscr L_q$ asigna $S_{90}(q)$ a cada generador
$[j,\mathcal F_6]$ y $S_{120}(q)$ al generador de frontera. Se obtiene
\begin{equation}
 \mathscr L_q(c_{12}^{+})=12S_{90}(q)-S_{120}(q).
 \label{iii:eq:barbero-weights}
\end{equation}
Los pesos $12$ y $-1$ son las multiplicidades orientadas de la cadena.
La orientación conjugada satisface $c_{12}^{-}=-c_{12}^{+}$, con los
mismos tipos de incidencia.

\begin{proposition}[Coeficiente singular del retorno fundamental]
\label{iii:prop:barbero-pole}
El coeficiente de $(1-q)^{-1}$ de
$\mathscr L_q(c_{12}^{+})$ es $1/24$.
\end{proposition}
\begin{proof}
Para cada entero $m>0$, la factorización
$1-q^{3m}=(1-q)\sum_{j=0}^{3m-1}q^j$ da
\[
 \lim_{q\uparrow1}(1-q)S_m(q)=\frac1{3m}.
\]
Aplicar esta identidad al funcional ya determinado en
\eqref{iii:eq:barbero-weights} produce
\begin{equation}
 \lim_{q\uparrow1}(1-q)\mathscr L_q(c_{12}^{+})
 =\frac{12}{270}-\frac1{360}=\frac1{24}.
 \label{iii:eq:barbero-pole}
\end{equation}
En la coordenada compleja $q-1$, el residuo es $-1/24$.
\end{proof}

La evaluación singular verifica una consecuencia de la incidencia y del
retorno trítico. El orden demostrativo queda fijado por
\eqref{iii:eq:barbero-flags}--\eqref{iii:eq:barbero-weights}; el
coeficiente singular se calcula después de construir la cadena.

